% from ega3-1-fr.tex
\oldpage[III]{111}
\begin{lemma}[2.5.3.1]
\label{III.2.5.3.1-fr}
Soient $A$ un corps, $r$ un entier $>0$, et $X=\mathbf P_A^r$~; on a
alors
\[
\chi_A(\mathcal O_X(n))=\binom{n+r}{r}
\]
pour tout $n\in\mathbb Z$.
\end{lemma}

\begin{proof}
En effet, pour $n\geq0$, on a
$\chi_A(\mathcal O_X(n))=\operatorname{long}H^0(X,\mathcal O_X(n))$, qui
est le nombre de monômes en les $T_i$ de degré total $n$, c'est-à-dire
$\binom{n+r}{r}$ (\hyperref[III.2.1.12-fr]{2.1.12}). Pour
$n\leq-r-1$, on a de même
\[
\chi_A(\mathcal O_X(n))=(-1)^r
\operatorname{long}H^r(X,\mathcal O_X(n))~;
\]
si $n=-r-h$, la dimension de $H^r(X,\mathcal O_X(n))$ sur $A$ est le
nombre des suites $(p_i)_{0\leq i\leq r}$ d'entiers $p_i>0$ tels que
$\sum_{i=0}^{r}p_i=r+h$ (\hyperref[III.2.1.12-fr]{2.1.12}), ou encore le
nombre des suites d'entiers $q_i\geq0$ $(0\leq i\leq r)$ tels que
$\sum_{i=0}^{r}q_i=h-1$~; c'est donc le nombre
\[
\binom{h+r-1}{r}=(-1)^r\binom{n+r}{r}.
\]
Enfin, pour $-r\leq n<0$, on a $\binom{n+r}{r}=0$ et d'autre part
$H^i(X,\mathcal O_X(n))=0$ pour tout $i\geq0$
(\hyperref[III.2.1.12-fr]{2.1.12}), ce qui prouve le lemme.
\end{proof}

\begin{corollary}[2.5.4]
\label{III.2.5.4-fr}
Soient $A$ un anneau local artinien, $S$ une $A$-algèbre graduée de type
fini engendrée par $S_1$, $M$ un $S$-module gradué de type fini,
$X=\operatorname{Proj}(S)$. On a alors
\[
\chi_A(\widetilde M(n))=\operatorname{long}M_n
\]
pour $n$ assez grand.
\end{corollary}

\begin{proof}
Cela résulte de ce que $M_n$ et $\Gamma(X,\widetilde M(n))$ sont
isomorphes pour $n$ assez grand
(\hyperref[III.2.3.1-fr]{2.3.1}).
\end{proof}

\subsection{Application : critères d'amplitude}
\label{subsection:III.2.6-fr}

\begin{proposition}[2.6.1]
\label{III.2.6.1-fr}
Soient $Y$ un préschéma noethérien, $f:X\to Y$ un morphisme propre,
$\mathcal L$ un $\mathcal O_X$-Module inversible. Les conditions
suivantes sont équivalentes~:
\begin{enumerate}
\item[a)] $\mathcal L$ est ample pour $f$.
\item[b)] Pour tout $\mathcal O_X$-Module cohérent $\mathcal F$, il
existe un entier $N$ tel que pour $n\geq N$, on ait
\[
R^qf_*(\mathcal F\otimes\mathcal L^{\otimes n})=0
\]
pour tout $q>0$.
\item[c)] Pour tout Idéal cohérent $\mathcal J$ de $\mathcal O_X$, il
existe un entier $N$ tel que pour $n\geq N$ on ait
\[
R^1f_*(\mathcal J\otimes\mathcal L^{\otimes n})=0.
\]
\end{enumerate}
\end{proposition}

\begin{proof}
On a vu que a) entraîne b) (\hyperref[III.2.2.1-fr]{2.2.1}, (ii)). Il
est trivial que b) entraîne c), et il reste à prouver que c) implique
a). On peut se borner au cas où $Y$ est affine (II, 4.6.4), et prouver
dans ce cas que $\mathcal L$ est ample~; il suffira de montrer que lorsque
$h$ parcourt l'ensemble des sections des $\mathcal L^{\otimes n}$
$(n>0)$ au-dessus de $X$, ceux des $X_h$ qui sont affines forment un
recouvrement de $X$ (II, 4.5.2). Pour cela, montrons que pour tout point
fermé $x$ de $X$ et tout voisinage ouvert affine $U$ de $x$, il existe un
$n$ et un $h\in\Gamma(X,\mathcal L^{\otimes n})$ tels que
$x\in X_h\subset U$~; $X_h$ sera nécessairement affine (I, 1.3.6) et la
réunion de ces $X_h$ sera un ouvert de $X$ contenant tous les points
fermés de $X$, et par suite $X$ lui-même puisque $X$ est noethérien
(I, 6.3.7 et 0\textsubscript{I}, 2.1.3). Soit $\mathcal J$ (resp.
$\mathcal J'$) le faisceau quasi-cohérent d'idéaux de $\mathcal O_X$
définissant le sous-préschéma fermé réduit de $X$ ayant pour espace
sous-jacent $X-U$ (resp. $(X-U)\cup\{x\}$) (I, 5.2.1)~; il est clair que
$\mathcal J$ et $\mathcal J'$ sont cohérents (I, 6.1.1), que
$\mathcal J'\subset\mathcal J$ et que
$\mathcal J''=\mathcal J/\mathcal J'$ est un $\mathcal O_X$-Module
cohérent (0\textsubscript{I}, 5.3.3) ayant pour support $\{x\}$ et tel
que $\mathcal J''_x=k(x)$. Comme $\mathcal L$ est localement libre, la
suite
\[
0\to\mathcal J'\otimes\mathcal L^{\otimes n}
\to\mathcal J\otimes\mathcal L^{\otimes n}
\to\mathcal J''\otimes\mathcal L^{\otimes n}\to0
\]
est exacte pour tout $n$, et par hypothèse il existe $n$ assez grand tel
que $H^1(X,\mathcal J'\otimes\mathcal L^{\otimes n})=0$~; la suite
exacte de cohomologie

